Over the past decade, Arabic Natural Language Processing (Arabic NLP) has transformed from a niche research area constrained by data scarcity into a vibrant, rapidly expanding field. Yet this growth has also revealed structural gaps, fragmented datasets, uneven dialect representation, and limited collaboration across institutions. As the community moves beyond the “resource-building” phase, the challenge is no longer just producing more data or larger models, but designing a sustainable ecosystem that reflects the linguistic and cultural realities of the Arab world. This keynote calls for reimagining Arabic NLP as an ecosystem rooted in representation, collaboration, and responsibility. Drawing on insights from large-scale projects such as MADAR, the Qatari Linguistic Map, and the LAILA Arabic Essay Scoring dataset, the talk will show how inclusive design, ethical data practices, and shared infrastructure can reshape how Arabic language technologies are developed and governed. It will highlight issues of bias, dialect homogenization, and access inequality, particularly in the era of generative AI, while outlining a vision for an Arabic NLP Commons, a framework for open data governance, equitable participation, and long-term community stewardship. Ultimately, the talk argues that success should be measured not only by technical achievements, but by how authentically it represents its speakers and empowers its researchers.
